% Unified manuscript. The bibliography is embedded at the end of this file.
% Compile with pdfLaTeX twice; no external bibliography or figure files are needed.
\documentclass[11pt]{article}
\usepackage[T1]{fontenc}
\usepackage{lmodern}
\usepackage[a4paper,margin=25mm]{geometry}
\usepackage{amsmath,amssymb,amsthm,mathtools}
\usepackage{microtype,booktabs,tabularx,array,enumitem}
\usepackage[numbers,sort&compress]{natbib}
\usepackage[hidelinks]{hyperref}
\usepackage{xurl}
\numberwithin{equation}{section}
\newtheorem{theorem}{Theorem}[section]
\newtheorem{proposition}[theorem]{Proposition}
\newtheorem{corollary}[theorem]{Corollary}
\newtheorem{lemma}[theorem]{Lemma}
\theoremstyle{definition}
\newtheorem{definition}[theorem]{Definition}
\theoremstyle{remark}
\newtheorem{remark}[theorem]{Remark}
\newcommand{\R}{\mathbb R}
\newcommand{\Tdet}{\mathsf T_{p,\mathrm{str}}^{\mathrm{det}}}
\newcommand{\Trand}{\mathsf T_{p,\mathrm{str}}^{\mathrm{rand}}}
\newcommand{\Tcert}{\mathsf T_{p,\mathrm{str}}^{\mathrm{det,cert}}}
\newcommand{\ap}{\alpha_p}
\newcommand{\logplus}{\log_{+}}
\DeclareMathOperator{\sym}{sym}
\newcolumntype{L}[1]{>{\raggedright\arraybackslash}p{#1}}
\newcolumntype{Y}{>{\raggedright\arraybackslash}X}
\setlength{\emergencystretch}{2em}
\setlength{\parskip}{0.12em}
\setlist{itemsep=0.25em,topsep=0.4em}
\clubpenalty=10000
\widowpenalty=10000
\displaywidowpenalty=10000
\title{Higher-Order Oracle Complexity of\\Strongly Monotone Variational Inequalities}
\author{Colombo Team\\[0.35em]
\small\href{mailto:colombo.2026@outlook.com}{\texttt{colombo.2026@outlook.com}}}
\date{}
\hypersetup{pdfauthor={Colombo Team},
pdftitle={Higher-Order Oracle Complexity of Strongly Monotone Variational Inequalities},
pdfsubject={Condition complexity, full-jet oracles, and tangent-residual certificates},
pdfkeywords={variational inequalities, strong monotonicity, higher-order methods, oracle complexity}}
\begin{document}
\maketitle
\begin{abstract}
We study the condition dependence of higher-order methods for strongly
monotone variational inequalities. For a fixed integer $p\ge2$, the operator
$F$ is $\mu$-strongly monotone on a known compact convex set of diameter at
most $D$, and $D^{p-1}F$ is $L_p$-Lipschitz. Each feasible query returns
$(F,DF,\ldots,D^{p-1}F)$; the target is a tangent residual of at most
$\varepsilon$. Write $\kappa_D=L_pD^{p-1}/\mu$ and
$\theta=\varepsilon/(\mu D)$. We give a complete proof that the optimal
polynomial condition exponent is $\alpha_p=2/(3p-1)$. A deterministic
algorithm returns an evaluated residual certificate within
$C_p(1+\kappa_D^{\alpha_p})\mathcal H^{6p-4}$ queries, where
$\mathcal H=1+\ln(2+\kappa_D+(\kappa_D+1)/\theta)$.
For sufficiently large $\kappa_D$ and $0<\theta\le1/16$, a matching
$c_p\kappa_D^{\alpha_p}$ lower bound holds in sufficiently high dimension
for arbitrary adaptive randomized algorithms with per-instance success
probability at least $2/3$. The proof specializes the extrapolated Halpern
hierarchy and flat-gated construction of the general monotone VI framework,
with explicit calibration to the prescribed strong-monotonicity modulus.
We also show that restarting a recent log-free monotone VI solver
sharpens the upper bound to
$O_p(1+\kappa_D^{\alpha_p}+\log_{+}(1/\theta))$.
Thus the condition dependence is sharp at fixed relative accuracy. One
scalar $p$-jet simulates one operator query, transferring the upper bounds
to strongly convex--strongly concave saddle problems.
\end{abstract}
\noindent\textbf{Keywords:} strongly monotone variational inequalities;
higher-order oracle complexity; tangent residual; randomized lower bounds;
Halpern iteration; saddle-point problems.
\par
\begin{samepage}
\noindent\textbf{AI Usage.}
Nearly the entire research pipeline for this paper was carried out by
Colombo, powered by GPT-5.6 Sol. 
\par
\end{samepage}

\section{Introduction}\label{sec:intro}
A variational inequality associated with an operator $F$ on a convex set
$\mathcal X$ asks for a point $x^\star\in\mathcal X$ satisfying
\[
\langle F(x^\star),x-x^\star\rangle\ge0
\qquad\text{for every }x\in\mathcal X.
\]
This formulation includes constrained optimization, convex--concave saddle
problems, and equilibrium systems that need not admit a scalar
potential~\cite{facchinei2003}. Strong monotonicity makes the solution
unique and turns a residual bound into a distance bound. The question
studied here is how the amount of higher-order information needed to find
that solution depends on the strong-monotonicity modulus.

Fix an order $p\ge2$. We assume that $F$ is $\mu$-strongly monotone,
$D^{p-1}F$ is $L_p$-Lipschitz, and $\mathcal X$ is known, compact, and
convex, with diameter at most $D$. A query at a feasible point returns the
complete jet $(F,DF,\ldots,D^{p-1}F)$. The requested output satisfies
\begin{equation}\label{eq:intro-residual}
r_{\mathrm{tan}}(x)
:=\operatorname{dist}\bigl(0,F(x)+N_{\mathcal X}(x)\bigr)
\le\varepsilon,
\end{equation}
where $N_{\mathcal X}$ is the outward normal cone. This criterion measures
stationarity at the returned point and remains meaningful when the
solution lies on the boundary. No common upper bound on $F$, $DF$, or the
other lower derivatives is imposed. The natural dimensionless parameters
are therefore
\[
\kappa_D=\frac{L_pD^{p-1}}{\mu},\qquad
\theta=\frac{\varepsilon}{\mu D}.
\]
Unlike a first-order condition number, $\kappa_D$ need not be at least one.

\paragraph{From tensor methods to accelerated resolvent solves.}
The extragradient and proximal-point frameworks provide classical methods
for monotone VIs~\cite{korpelevich1976,rockafellar1976,nemirovski2004}.
Newton and higher-order extensions replace the operator by a regularized
Taylor model before applying an extragradient
correction~\cite{monteiro2012,bullins2020,adil2022,lin2024}.
Restarting such methods gives the strong-VI condition exponent
$2/(p+1)$~\cite{ostroukhov2020}. More recent methods exploit the structure
of successive resolvent problems. Higher-order Halpern iteration yields
the general monotone exponent $1/p$~\cite{chen2026halpern}, while
extrapolated Halpern methods~\cite{zhang2026,zhangchen2026} and anchored
extra-proximal methods~\cite{jiangyoon2026} reach $2/(3p-1)$.
These developments make it necessary to distinguish a general monotone
accuracy bound from a condition bound with a prescribed $\mu$.

\paragraph{The condition-complexity statement.}
Put
\[
\ap=\frac{2}{3p-1},\qquad
\mathcal H=1+\ln\!\left(2+\kappa_D+
\frac{\kappa_D+1}{\theta}\right).
\]
Our complete, self-contained construction proves
\begin{equation}\label{eq:intro-frontier}
c_p\kappa_D^{\ap}
\le\Trand\le\Tdet\le\Tcert
\le C_p(1+\kappa_D^{\ap})\mathcal H^{6p-4}
\end{equation}
whenever $\kappa_D\ge\kappa_{0,p}$ and $0<\theta\le1/16$.
Here the randomized complexity allows arbitrary adaptive queries and
arbitrary feasible outputs, with success probability at least $2/3$ on
each instance. The deterministic certificate complexity requires the
stronger output $(x,n,F(x))$, with $n\in N_{\mathcal X}(x)$ and
$\|F(x)+n\|\le\varepsilon$, evaluated at a charged query. The upper
bound holds in every finite dimension and for all positive public
parameters. The lower bound is realized on a Euclidean ball in an
explicitly specified dimension that grows with the condition parameter.

The sharpest upper comparison uses the recent log-free monotone
solver of Zhang, Chen, Pan, Liu, Yang, and
Zhang~\cite[Theorem~5.1]{zhangchen2026}. A short restart argument gives
\begin{equation}\label{eq:intro-sharper}
\Tcert\le C_p\left(1+\kappa_D^{\ap}
+\logplus\frac1\theta\right),
\qquad \logplus t=\max\{0,\ln t\}.
\end{equation}
We prove this consequence separately in Proposition~\ref{prop:literature-restart}
and identify its external algorithmic input explicitly. Combining it with
our randomized lower bound removes the logarithmic condition loss at any
fixed $\theta\in(0,1/16]$. Neither bound determines the optimal dependence
on arbitrarily small $\theta$.

\paragraph{Contribution and attribution.}
The exponent $2/(3p-1)$, the extrapolated Halpern hierarchy, and the
flat-gated hard family already appear in the general monotone VI analysis
of Zhang et al.~\cite{zhang2026}. This paper gives their prescribed-modulus
specialization with a complete proof: the lower family satisfies the
public $\mu$ and $L_p$ simultaneously; the upper algorithm accepts only
evaluated certificates and charges all nested calls; and the randomized
lower bound fixes the instance before the algorithm's seed is realized.
The scalar saddle reduction preserves both the residual and the number
of calls. The sharper estimate~\eqref{eq:intro-sharper} is a consequence
of~\cite{zhangchen2026}, rather than a new monotone VI algorithm.

\paragraph{Why the exponent has this value.}
The upper and lower arguments expose complementary mechanisms. On the
upper side, firm nonexpansiveness controls the squared second differences
of a Halpern orbit over a block of $k$ iterations by
$O(R^2\log^2(k+2)/k^2)$. A projected linear predictor therefore makes the
aggregate cost of successive regularized solves small. If the inner
strong solver has exponent $\gamma$, the outer cost is, up to logarithms,
\[
K+\left(\frac{L_pR^p}{\varepsilon K}\right)^\gamma.
\]
Balancing the two terms gives $\gamma/(1+\gamma)$, and certified restart
returns a strong solver with that exponent. Starting from $2/(p+1)$,
exactly $p-1$ admissible layers give $2/(3p-1)$.
On the lower side, a contractive chain with $m$ hidden directions has
strong-monotonicity modulus proportional to
$L_pD^{p-1}m^{-(3p-1)/2}$. Taking
$m\asymp_p\kappa_D^{2/(3p-1)}$ preserves the prescribed modulus while
keeping the residual above $\mu D/8$ until the last direction is revealed.

\paragraph{Organization.}
Section~\ref{sec:related} aligns the relevant literature by oracle and
error criterion. Sections~\ref{sec:setup} and~\ref{sec:results} specify
the model and theorems. Section~\ref{sec:upper} proves the self-contained
upper bound; Section~\ref{sec:lower} proves the randomized lower bound.
Section~\ref{sec:saddle} gives the scalar saddle reduction, and
Section~\ref{sec:discussion} states the remaining limitations.

\section{Related Work and Rate Comparisons}\label{sec:related}
All orders in this section follow the same convention: an order-$p$
operator oracle returns derivatives through $D^{p-1}F$, and the
corresponding scalar oracle returns derivatives through order $p$.
Constants denoted by $O_p$ depend only on the fixed order.

\paragraph{Error criteria.}
For a monotone operator, define the Minty and Stampacchia gaps by
\[
g_{\mathrm M}(x)=\sup_{u\in\mathcal X}\langle F(u),x-u\rangle,
\qquad
g_{\mathrm S}(x)=\sup_{u\in\mathcal X}\langle F(x),x-u\rangle.
\]
Monotonicity and the normal-cone inequality imply
\begin{equation}\label{eq:gap-comparison}
g_{\mathrm M}(x)\le g_{\mathrm S}(x)\le D r_{\mathrm{tan}}(x).
\end{equation}
Indeed, for $n\in N_{\mathcal X}(x)$,
$\langle F(x),x-u\rangle\le\langle F(x)+n,x-u\rangle$.
A tangent certificate therefore controls both gaps, but a small gap does
not generally control the tangent residual at the same point. For example,
on $[0,1]$, $F(x)=1+x$ is strongly monotone and has solution zero;
$g_{\mathrm S}(x)=x(1+x)\to0$ as $x\downarrow0$, whereas
$r_{\mathrm{tan}}(x)=1+x$ for every interior $x$. This distinction is
essential when comparing published rates.

\paragraph{Classical and higher-order methods.}
Extragradient, Mirror-Prox, and dual extrapolation provide the classical
first-order foundations~\cite{korpelevich1976,nemirovski2004,nesterov2007}.
The Newton proximal extragradient method~\cite{monteiro2012}, higher-order
Mirror-Prox~\cite{bullins2020}, and subsequent tensor
methods~\cite{adil2022,lin2024} achieve the characteristic gap exponent
$2/(p+1)$. The lower bounds in the latter two works use restrictions on
the generated iterates; they do not cover every adaptive full-jet policy.
Ostroukhov et al.~\cite{ostroukhov2020} develop restarted tensor methods
for strongly monotone VIs and strongly convex--strongly concave saddle
problems. Their condition exponent provides the historical baseline.

\paragraph{Anchoring and recent general monotone results.}
Halpern-based methods connect fixed-point residuals to strong VI
solutions~\cite{diakonikolas2020}; tangent residuals also underpin
last-iterate guarantees~\cite{cai2022}. Chen et
al.~\cite{chen2026halpern} combine Halpern iteration with tensor solves,
obtaining exponent $1/p$ for a normalized proximal residual. Their
residual comparison controls the tangent residual at the associated
exact resolvent point, which need not be the reported point. Zhang et
al.~\cite{zhang2026} use extrapolated warm starts and a finite hierarchy
to obtain exponent $\ap$ and an unrestricted randomized lower bound.
The recent method of Zhang, Chen, Pan, Liu, Yang, and
Zhang~\cite{zhangchen2026} attains the same upper exponent without
logarithmic losses. Jiang and Yoon~\cite{jiangyoon2026} obtain it, up to
a logarithm, for composite monotone inclusions, and prove an unrestricted
deterministic lower bound. Their single-valued operator is globally
smooth, and Taylor centers need not belong to the constrained domain;
their oracle model therefore differs from our feasible-query model.
Table~\ref{tab:general-comparison} records the main distinctions.

\begin{table}[t]
\centering\small
\setlength{\tabcolsep}{4pt}
\renewcommand{\arraystretch}{1.18}
\begin{tabularx}{\textwidth}{@{}L{0.25\textwidth}L{0.17\textwidth}L{0.20\textwidth}Y@{}}
\toprule
Work / setting & Criterion & Upper queries & Lower-bound scope\\
\midrule
Tensor VI methods\newline\cite{adil2022,lin2024}
& Minty gap & $O_p(1+Q_{\rm gap}^{2/(p+1)})$
& The cited matching lower bounds impose tensor or span restrictions.\\
Chen et al.\ \cite{chen2026halpern}
& Proximal residual
& $\widetilde O_p(1+Q_{\rm res}^{1/p})$
& Upper-bound result for general monotone VIs.\\
Zhang et al.\ \cite{zhang2026}
& Evaluated tangent certificate
& $\widetilde O_p(1+Q_{\rm res}^{\ap})$
& $\Omega_p(Q_{\rm res}^{\ap})$ for arbitrary adaptive randomized policies.\\
Zhang--Chen--Pan et al.\ \cite{zhangchen2026}
& Evaluated tangent certificate
& $O_p(1+Q_{\rm res}^{\ap})$
& Upper bound; comparisons with lower bounds inherit their stated scope.\\
Jiang--Yoon\ \cite{jiangyoon2026}
& Tangent residual; global smoothness
& $O_p((1+Q_{\rm res}^{\ap})\ell_Q)$
& Composite inclusions; unrestricted deterministic lower bound;
queries need not be feasible.\\
\bottomrule
\end{tabularx}
\caption{General monotone benchmarks, in a common order convention.
Here $Q_{\rm res}=L_pD^p/\varepsilon$, with the residual criterion
specified separately in each row, and
$Q_{\rm gap}=L_pD^{p+1}/\varepsilon_G$, and
$\ell_Q=\ln(2+Q_{\rm res})$. These are different accuracy
parameters. The table uses the domain diameter $D$ as a valid distance
bound for bounded VIs; the composite-inclusion row uses a supplied
distance bound in its place. Some sources allow the sharper
initial-distance bound $R$. All unrestricted lower bounds require
sufficiently high dimension.}
\label{tab:general-comparison}
\end{table}

\paragraph{Strong monotonicity and the present comparison.}
The general monotone rate $Q_{\rm res}^{\ap}$ becomes
$(\kappa_D/\theta)^{\ap}$ if applied directly. Restart uses
$\mu\|x-x^\star\|\le r_{\mathrm{tan}}(x)$ to improve the accuracy
dependence. Our self-contained proof yields the explicit factor
$\mathcal H^{6p-4}$ and an initial-distance refinement.
Proposition~\ref{prop:literature-restart} gives the stronger consequence
of the log-free solver. Thus the contribution is a calibrated
strong-monotonicity statement with full randomized lower-bound
quantifiers, not a claim to improve the best known general monotone
upper exponent or its logarithmic factors.

\paragraph{Saddle problems and hidden derivatives.}
Chen et al.~\cite{chen2025} obtain second-order gap complexity
$\widetilde O(\varepsilon_G^{-4/7})$ for convex--concave minimax
problems; their higher-order development~\cite{chen2026saddle} gives
tangent-residual exponent $4/(3p+1)$ and a lower exponent $\ap$ under
its tensor-algorithm restriction. These are fixed-geometry accuracy
comparisons; the higher-order upper analysis also uses a first-derivative
smoothness bound in logarithmic factors. General VI upper bounds transfer to
scalar saddle problems through the signed gradient. A general VI lower
bound does not transfer in the reverse direction: the operator must also
have the required scalar integrability structure.
The exact flat gates used here follow the information-hiding approach
in~\cite{zhang2026}, related to the smooth contractive fixed-point
construction of Jang and Ryu~\cite{jang2026}. We verify explicitly that
the gates hide every tensor in a reply and that boundary normals cannot
cancel the lower-bound residual.


\section{Problem Formulation}\label{sec:setup}
We count information acquired from the unknown operator. All other
computations use the public domain or information already returned by the
oracle. This distinction allows arbitrary convex domains while keeping
the complexity statement independent of how tensors and constraints are
represented.

\subsection{Public data and the operator class}
Fix an integer $p\ge2$. The public data are a finite dimension $d\ge1$,
positive numbers $L_p,\mu,D,\varepsilon$, a known nonempty compact convex
set $\mathcal X\subseteq\R^d$ with $\operatorname{diam}(\mathcal X)\le D$,
and a start $x_0\in\mathcal X$. All vector norms are Euclidean. For a
vector-valued $j$-linear map $A$, use the norm
\[
\|A\|_{\mathrm{op}}
=\sup_{\|u_1\|,\ldots,\|u_j\|\le1}
\|A[u_1,\ldots,u_j]\|.
\]
The unknown operator $F$ is defined and $C^{p-1}$ on an open neighborhood
of $\mathcal X$. For all $x,y\in\mathcal X$, it satisfies
\begin{align}
\langle F(x)-F(y),x-y\rangle
&\ge\mu\|x-y\|^2,\label{eq:strong-monotonicity}\\
\|D^{p-1}F(x)-D^{p-1}F(y)\|_{\mathrm{op}}
&\le L_p\|x-y\|.\label{eq:operator-smoothness}
\end{align}
We denote this class by
$\mathfrak F^{\mathrm{str}}_{p,L_p,\mu}(\mathcal X)$.
The constants are one-sided promises: $\mu$ is a lower bound on the
strong-monotonicity modulus, and $L_p$ is an upper bound on the highest
derivative's Lipschitz modulus. We do not require a continuous $D^pF$, a
scalar potential, or uniform bounds on the lower derivatives. The class
is nonempty for every positive parameter tuple, since
$F(x)=\mu x+b$ satisfies both promises for any $b\in\R^d$.

The solution $x^\star$ is characterized by
\begin{equation}\label{eq:vi-solution}
\langle F(x^\star),x-x^\star\rangle\ge0
\qquad(x\in\mathcal X).
\end{equation}
Existence follows from the fixed-point theorem applied to the continuous
map $x\mapsto\Pi_{\mathcal X}(x-F(x))$. Adding the inequalities for two
solutions and using~\eqref{eq:strong-monotonicity} proves uniqueness. Define the outward
normal cone and tangent residual by
\begin{align}
N_{\mathcal X}(x)
&=\{n:\langle n,u-x\rangle\le0\text{ for every }u\in\mathcal X\},
\label{eq:normal-cone}\\
r_{\mathrm{tan}}(x)
&=\min_{n\in N_{\mathcal X}(x)}\|F(x)+n\|.
\label{eq:tangent-residual}
\end{align}
The minimum exists because the normal cone is nonempty and closed.
At a solution, $-F(x^\star)\in N_{\mathcal X}(x^\star)$, including when
$F(x^\star)$ is nonzero.

\begin{lemma}[Residual controls distance]\label{lem:error-bound}
For every $x\in\mathcal X$ and $n\in N_{\mathcal X}(x)$,
\begin{equation}\label{eq:error-distance}
\mu\|x-x^\star\|\le r_{\mathrm{tan}}(x)\le\|F(x)+n\|.
\end{equation}
\end{lemma}
\begin{proof}
The normal-cone inequalities at $x$ and $x^\star$ imply
\[
\langle F(x)+n,x-x^\star\rangle
\ge\langle F(x)-F(x^\star),x-x^\star\rangle
\ge\mu\|x-x^\star\|^2.
\]
Apply Cauchy--Schwarz and minimize over $n$.
\end{proof}

\subsection{Oracle, policies, and evaluated certificates}
One charged query at a feasible point $x$ returns
\begin{equation}\label{eq:full-jet}
\mathcal O_p^F(x)=(F(x),DF(x),\ldots,D^{p-1}F(x)).
\end{equation}
All returned tensors are exact. Every call to this oracle counts,
including calls inside regularized problems, unsuccessful radius trials,
restarts, and the final certificate check. Projections onto the known
domain and exact solution of a model formed solely from returned jets
are internal computations. An evaluation of the unknown operator's
resolvent is implemented through charged queries; it is not supplied as
an additional oracle.

Policies may use the complete transcript, arbitrary Borel measurable
internal computations, and any feasible output, whether queried or not.
A randomized policy uses an instance-independent seed on a standard
Borel space. Query feasibility and the query cap hold on every sample
path, and each fixed seed defines a legal deterministic policy.
Its success probability must be at least $2/3$ on every fixed instance.
Thus an adversarial instance may depend on the policy but must be fixed
before the realized seed. Policy rules are taken to be total Borel maps
on transcripts; outside realizable transcripts they can be extended by
feasible default queries and outputs. Equivalently, projecting such
default points onto $\mathcal X$ enforces feasibility without changing
any legal run.

For the public tuple $(p,L_p,\mu,D,\varepsilon)$, let $\Tdet$ and
$\Trand$ be the least nonnegative integer query caps for deterministic
and randomized point output. The worst case ranges over every finite
dimension, every legal public domain and start, and every compatible
operator. Let $\Tcert$ be the deterministic complexity for returning
an evaluated triple
\begin{equation}\label{eq:evaluated-certificate}
(\widehat x,n,F(\widehat x)),\qquad
n\in N_{\mathcal X}(\widehat x),\qquad
\|F(\widehat x)+n\|\le\varepsilon,
\end{equation}
where $F(\widehat x)$ was obtained at a charged query at that same point.
If no finite uniform cap exists, the corresponding complexity is
$+\infty$. These definitions immediately imply
\begin{equation}\label{eq:complexity-order}
\Trand\le\Tdet\le\Tcert.
\end{equation}
For lower bounds, an early stop can be padded by ignored feasible
queries while preserving its saved output.

\subsection{Normalization}
Throughout the paper,
\begin{equation}\label{eq:normalized-parameters}
\kappa_D=\frac{L_pD^{p-1}}{\mu},\qquad
\theta=\frac{\varepsilon}{\mu D},\qquad
\ap=\frac{2}{3p-1},
\end{equation}
and
\begin{equation}\label{eq:public-log}
\mathcal H(\kappa_D,\theta)
=1+\ln\!\left(2+\kappa_D+\frac{\kappa_D+1}{\theta}\right).
\end{equation}
The notation $c_p,C_p$ denotes positive finite constants depending only
on $p$; their values may change between displays. Specified constants,
such as $c_p^{\mathrm{sc}}$ in the lower-bound construction, retain their
defined values. The diameter bound $D$ is unrelated to the derivative
notation $D^jF$.

\section{Main Results}\label{sec:results}
We first state the bounds proved entirely within this paper. We then
record the stronger upper estimate obtained by importing the recent
log-free monotone solver. This separates the two proof dependencies.

\begin{theorem}[Self-contained certified upper bound]\label{thm:upper}
For each fixed $p\ge2$, there is a total deterministic measurable
algorithm that, for every positive $L_p,\mu,D,\varepsilon$, every finite
dimension, and every legal instance, returns~\eqref{eq:evaluated-certificate} within
\begin{equation}\label{eq:hierarchy-upper}
\Tcert(\varepsilon;L_p,\mu,D)
\le C_p(1+\kappa_D^{\ap})\mathcal H(\kappa_D,\theta)^{6p-4}
\end{equation}
queries. More precisely, on an instance with
$r_0=\|x_0-x^\star\|$, its cost is at most
\begin{equation}\label{eq:distance-sensitive-upper}
C_p\left[1+\left(\frac{L_pr_0^{p-1}}{\mu}\right)^{\ap}\right]
\mathcal H(\kappa_D,\theta)^{6p-4}.
\end{equation}
Neither $r_0$ nor $x^\star$ is supplied to the algorithm.
\end{theorem}

Section~\ref{sec:upper} constructs the algorithm by a finite hierarchy
starting from tensor extragradient. Public query caps make unsuccessful
radius guesses safe, and an evaluated certificate is the only acceptance
test. In particular, the initial-distance refinement holds even when
$r_0=0$.

\begin{theorem}[Randomized lower bound with prescribed modulus]
\label{thm:lower}
For each fixed $p\ge2$, there are $c_p,\kappa_{0,p}>0$ such that, whenever
\begin{equation}\label{eq:lower-regime}
\kappa_D\ge\kappa_{0,p},\qquad
0<\varepsilon\le\frac{\mu D}{16},
\end{equation}
every randomized adaptive policy that succeeds with probability at least
$2/3$ on every legal instance has uniform query cap at least
\begin{equation}\label{eq:randomized-lower}
\Trand(\varepsilon;L_p,\mu,D)
\ge c_p\kappa_D^{\ap}.
\end{equation}
The lower bound permits arbitrary feasible full-jet queries and arbitrary
feasible outputs. The hard instance is fixed before the policy's seed.
\end{theorem}

The proof uses the public ball $B_2^d(D/2)$ and the public start zero.
With the constant $c_p^{\mathrm{sc}}$ specified
in~\eqref{eq:lower-constants}, set
\[
m=\left\lfloor(c_p^{\mathrm{sc}}\kappa_D)^{\ap}\right\rfloor.
\]
Any integer dimension satisfying
\[
d\ge m+\left\lceil2048m^3\ln\bigl(8(m^2+1)\bigr)\right\rceil
\]
suffices. Every policy using at most $m-1$ calls fails on one fixed legal
operator with probability at least $3/4$; its residual is then at least
$2\varepsilon$. The dimension requirement is sufficient, not asserted
to be optimal.

\subsection{A sharper upper bound from a recent monotone solver}
The next proposition uses an external theorem as an explicit input.
It improves the logarithmic factors of Theorem~\ref{thm:upper}, while
the proof in Section~\ref{sec:upper} remains independent of that input.

\begin{proposition}[Restart of the log-free monotone solver]
\label{prop:literature-restart}
The monotone VI algorithm of Zhang, Chen, Pan, Liu, Yang, and
Zhang~\cite[Theorem~5.1]{zhangchen2026}, followed by certified restart,
gives, for all positive public parameters,
\begin{equation}\label{eq:sharper-upper}
\Tcert(\varepsilon;L_p,\mu,D)
\le C_p\left(1+\kappa_D^{\ap}+
\logplus\frac{\mu D}{\varepsilon}\right).
\end{equation}
\end{proposition}
\begin{proof}
We use the cited theorem in the following radius form. Given a monotone
operator with the smoothness promise~\eqref{eq:operator-smoothness}, a feasible start
$w$, and a known $R>0$ with $\|w-z^\star\|\le R$ for a solution
$z^\star$, its algorithm returns an evaluated tangent certificate of
norm at most $\delta$ in
\begin{equation}\label{eq:imported-radius-solver}
C_p\left[1+\left(\frac{L_pR^p}{\delta}\right)^{\ap}\right]
\end{equation}
feasible full-jet queries. In the source, the symbol $D$ denotes the
initial distance. The proof uses this distance only through upper
bounds: replacing it by $R$ in the initial step size
$1/(L_pR^{p-1})$ and inner tolerance $R/T^{3/2}$ preserves those
bounds. The source's projected warm starts are feasible, and its inner
acceptance test evaluates the operator at the certified point. Thus its
output and query convention agree with ours. The affine-hull reduction
in Section~\ref{sec:upper} handles lower-dimensional domains.

Set $R_j=D2^{-j}$ and
\[
\delta_j=\max\{\varepsilon,\mu R_j/2\}.
\]
Start stage zero at $x_0$. At stage $j$, call~\eqref{eq:imported-radius-solver}
on the original domain with radius $R_j$ and tolerance $\delta_j$.
If $\delta_j=\varepsilon$, return its certificate. Otherwise,
Lemma~\ref{lem:error-bound} places the returned point within $R_j/2$
of $x^\star$, so it is a valid start for stage $j+1$.
The last stage is
\[
J=\max\!\left\{0,\left\lceil
\log_2\frac{\mu D}{2\varepsilon}\right\rceil\right\}.
\]
Since $\delta_j\ge\mu R_j/2$, the total cost is at most
\begin{align*}
C_p\sum_{j=0}^{J}
\left[1+\left(\frac{2L_pR_j^{p-1}}{\mu}\right)^{\ap}\right]
&\le C_p\left[J+1+
\frac{(2\kappa_D)^{\ap}}{1-2^{-(p-1)\ap}}\right]\\
&\le C_p\left(1+\kappa_D^{\ap}+
\logplus\frac{\mu D}{\varepsilon}\right).
\end{align*}
All radii and tolerances are public. No new constraint is introduced
during a restart, so every normal in the output is a normal to the
original domain. The cited finite call bounds may be installed as public
caps, and every final evaluation is included in the count.
\end{proof}

\begin{corollary}[Sharp condition dependence at fixed relative accuracy]
\label{cor:frontier}
In the regime~\eqref{eq:lower-regime},
\begin{equation}\label{eq:combined-frontier}
c_p\kappa_D^{\ap}
\le\Trand\le\Tdet\le\Tcert
\le C_p\left(1+\kappa_D^{\ap}+\logplus\frac1\theta\right).
\end{equation}
Consequently, for any fixed $\theta\in(0,1/16]$ and
$\kappa_D\to\infty$,
\[
\Trand\asymp_p\Tdet\asymp_p\Tcert\asymp_p\kappa_D^{\ap},
\]
where the asymptotic threshold may depend on the fixed $\theta$.
Using only the self-contained upper proof gives
$\widetilde\Theta_p(\kappa_D^{\ap})$ instead.
\end{corollary}
\begin{proof}
Combine Theorem~\ref{thm:lower}, Proposition~\ref{prop:literature-restart},
and~\eqref{eq:complexity-order}. At fixed $\theta$, the additive accuracy term is
constant. The self-contained alternative follows from
Theorem~\ref{thm:upper} and
$\mathcal H(\kappa_D,\theta)=O(\ln\kappa_D)$.
\end{proof}

\begin{table}[t]
\centering\small
\setlength{\tabcolsep}{4pt}
\renewcommand{\arraystretch}{1.18}
\begin{tabularx}{\textwidth}{@{}L{0.25\textwidth}L{0.34\textwidth}Y@{}}
\toprule
Result & Query bound & Criterion and scope\\
\midrule
Ostroukhov et al.\ \cite{ostroukhov2020}
& $O_p\!\left(\kappa_R^{2/(p+1)}
\ln\frac{\mu R^2}{\varepsilon_G}\right)$
& Source's gap guarantee in its small-error regime; unconstrained setting,
$\kappa_R=L_pR^{p-1}/\mu$.\\
Certified tensor baseline\newline
Lemma~\ref{lem:strong-base}
& $O_p((1+\kappa_D^{2/(p+1)})\mathcal H^2)$
& Evaluated tangent certificate; current model.\\
Self-contained hierarchy\newline
Theorem~\ref{thm:upper}
& $O_p((1+\kappa_D^{\ap})\mathcal H^{6p-4})$
& All positive parameters; all finite dimensions.\\
Restart of\ \cite{zhangchen2026}\newline
Proposition~\ref{prop:literature-restart}
& $O_p(1+\kappa_D^{\ap}+\logplus(1/\theta))$
& Same evaluated certificate and feasible-query model.\\
Calibrated randomized lower bound\newline
Theorem~\ref{thm:lower}
& $\Omega_p(\kappa_D^{\ap})$
& $\kappa_D\ge\kappa_{0,p}$, $\theta\le1/16$;
arbitrary adaptive policies; high dimension.\\
\bottomrule
\end{tabularx}
\caption{Strong-monotonicity comparison. Only the first row uses a
different error criterion and domain. All other rows use
the exact feasible full-jet model of Section~\ref{sec:setup}.
The sharper upper row explicitly depends on the algorithm in the cited
paper. The exponent $\ap$ equals $2/5$, $1/4$, and $2/11$ for
$p=2,3,4$, respectively.}
\label{tab:comparison}
\end{table}


\section{A Self-Contained Certified Upper Bound}\label{sec:upper}
We prove Theorem~\ref{thm:upper} through a finite hierarchy of solvers.
A regularized Taylor step first gives a distance-sensitive strongly
monotone solver. A restart procedure and an extrapolated Halpern procedure
then improve its exponent. Throughout the construction, a successful call
returns a normal vector and the operator value evaluated at the same
point. Every nested call is charged to the original oracle.

\subsection{Affine-hull reduction and Taylor certificates}

We first work in the known affine hull of the feasible set. Write
\(\operatorname{aff}\mathcal X=x_c+U\mathbb R^m\), with
\(U^\top U=I_m\), and set
\[
\widetilde{\mathcal X}=\{u:x_c+Uu\in\mathcal X\},
\qquad
\widetilde F(u)=U^\top F(x_c+Uu).
\]
One ambient jet supplies the entire reduced jet, with unchanged
smoothness and strong-monotonicity promises. A reduced certificate lifts
without increasing its norm: for \(x=x_c+Uu\), define
\[
n=U\widetilde n-(I-UU^\top)F(x).
\]
Then \(n\in N_{\mathcal X}(x)\) and
\(F(x)+n=U(\widetilde F(u)+\widetilde n)\).
If \(m=0\), one query at the unique feasible point and the normal
\(-F(x)\) provide an exact certificate. We may therefore assume below
that \(\mathcal X\) has nonempty interior in its working space.

\begin{lemma}[Regularized Taylor certificate]\label{lem:taylor-certificate}
Let \(G\) be monotone on \(\mathcal X\), with
\(\operatorname{Lip}(D^{p-1}G)\le L_p\). For \(x\in\mathcal X\), define
\[
\overline G_x(y)=\sum_{j=0}^{p-1}\frac{1}{j!}D^jG(x)[y-x]^j,
\qquad \beta=\frac{2L_p}{(p-1)!}.
\]
The model variational inequality
\begin{equation}\label{eq:taylor-model}
0\in\overline G_x(y)+\beta\|y-x\|^{p-1}(y-x)
       +N_{\mathcal X}(y)
\end{equation}
has a unique solution. If \(z^\star\) solves the VI for \(G\) and
\(\|x-z^\star\|\le r\), the model solution and its associated normal
provide an evaluated certificate of norm at most \(A_pL_pr^p\) using
at most two oracle calls, where
\(A_p=(2p+1)3^p/p!\).
\end{lemma}

\begin{proof}
The integral Taylor remainder gives
\begin{equation}\label{eq:taylor-remainder}
\|G(y)-\overline G_x(y)\|
\le\frac{L_p}{p!}\|y-x\|^p.
\end{equation}
Applying the same estimate to \(DG\), one order lower, gives
\[
\|DG(y)-D\overline G_x(y)\|_{\mathrm{op}}
\le\frac{L_p}{(p-1)!}\|y-x\|^{p-1}.
\]
These estimates require Lipschitz continuity of \(D^{p-1}G\), rather
than the existence of a continuous \(D^pG\).
Monotonicity implies that the symmetric part of \(DG\) is positive
semidefinite in the interior; continuity extends this conclusion to the
boundary. The derivative of \(s\mapsto\beta\|s\|^{p-1}s\) is bounded
below by \(\beta\|s\|^{p-1}I\). Consequently, the symmetric Jacobian of
the model operator is bounded below by
\(L_p\|y-x\|^{p-1}I/(p-1)!\).
Its integral along any nonconstant feasible segment is strictly positive.
The model is therefore strictly monotone. Continuity and compactness give
existence, and strict monotonicity gives uniqueness.

The model solution supplies the exact normal
\begin{equation}\label{eq:taylor-normal}
n=-\overline G_x(y)-\beta\|y-x\|^{p-1}(y-x)
\in N_{\mathcal X}(y).
\end{equation}
Put \(s=y-x\), \(E=G(y)-\overline G_x(y)\), and
\(v=G(y)+n=E-\beta\|s\|^{p-1}s\).
Monotonicity of \(G+N_{\mathcal X}\) gives
\(\langle v,y-z^\star\rangle\ge0\). If \(s\ne0\),
\eqref{eq:taylor-remainder} and \(\|x-z^\star\|\le r\) imply
\[
\beta(\|s\|-r)\le\frac{L_p}{p!}(\|s\|+r).
\]
Since \(\beta=2pL_p/p!\), it follows that
\(\|s\|\le(2p+1)r/(2p-1)\le3r\). Thus
\begin{equation}\label{eq:taylor-certificate}
\|y-x\|\le3r,
\qquad
\|G(y)+n\|\le A_pL_pr^p,
\qquad A_p=\frac{(2p+1)3^p}{p!}.
\end{equation}
A jet at \(x\) forms the model, and a query at \(y\) evaluates the
certificate. If \(y=x\), the first query already evaluates an exact
certificate.
\end{proof}

\subsection{A distance-sensitive base solver}

For positive \(\nu,\delta\), define the public logarithmic factor
\[
\mathcal L_S(\nu,\delta)
=1+\ln\!\left(2+\frac{L_pD^{p-1}}{\nu}
                +\frac{L_pD^p}{\delta}
                +\frac{\nu D}{\delta}\right).
\]
A \emph{strong solver with exponents \((\gamma,b)\)} is a deterministic
routine for a \(\nu\)-strongly monotone operator \(G\), a supplied start
\(w\), and a tolerance \(\delta\). It returns an evaluated certificate
of norm at most \(\delta\), with cost
\[
C_S\left[1+\left(\frac{L_pr^{p-1}}{\nu}\right)^\gamma\right]
\mathcal L_S(\nu,\delta)^b,
\qquad r=\|w-z^\star\|,
\]
without receiving \(r\) or \(z^\star\). The constant \(C_S\) is fixed
when the routine is constructed. All exponents and constants used below
are functions of the fixed order \(p\).

\begin{lemma}[Base strongly monotone solver]\label{lem:strong-base}
There is a strong solver with
\(\gamma_0=2/(p+1)\) and \(b_0=2\).
\end{lemma}

\begin{proof}
First suppose a positive radius \(r\) is a valid bound on the current
distance to \(z^\star\). Form the model solution from
Lemma~\ref{lem:taylor-certificate}. If \(s=y-x\ne0\), set
\[
\lambda=(\beta\|s\|^{p-1})^{-1},
\qquad x^+=\Pi_{\mathcal X}(x-\lambda G(y)).
\]
This iteration uses at most two original-operator queries. To verify its
energy inequality, write \(d=x^+-x\). The projection inequality and
\(\langle n,y-x^+\rangle\ge0\) give, for \(u\in\mathcal X\),
\begin{align*}
2\lambda\langle G(y),y-u\rangle
&\le\|x-u\|^2-\|x^+-u\|^2-\|d\|^2
       +2\lambda\langle G(y)+n,y-x^+\rangle\\
&=\|x-u\|^2-\|x^+-u\|^2-\|d-s\|^2-\|s\|^2
       -2\langle\lambda E,d-s\rangle.
\end{align*}
Here \(\lambda(G(y)+n)=\lambda E-s\), and
\(\|\lambda E\|\le\|s\|/(2p)\). Completing the square proves
\begin{equation}\label{eq:extragradient-descent}
2\lambda\langle G(y),y-u\rangle
\le\|x-u\|^2-\|x^+-u\|^2-\frac34\|s\|^2.
\end{equation}
For \(u=z^\star\), strong monotonicity and the solution VI imply
\(\langle G(y),y-z^\star\rangle\ge\nu\|y-z^\star\|^2\).
Summing \(T\) iterations therefore gives
\begin{equation}\label{eq:extragradient-sum}
2\nu\sum_{i<T}\lambda_i\|y_i-z^\star\|^2
+\frac34\sum_{i<T}\|s_i\|^2\le r^2.
\end{equation}
If some \(s_i=0\), that iteration already returns an exact certificate.
Otherwise, convexity of \(t\mapsto t^{-(p-1)/2}\) yields
\begin{equation}\label{eq:stepsize-sum}
\Lambda_T:=\sum_{i<T}\lambda_i
\ge a_p\frac{T^{(p+1)/2}}{L_pr^{p-1}},
\qquad
 a_p=\frac{(p-1)!}{2}\left(\frac34\right)^{(p-1)/2}.
\end{equation}
The feasible weighted mean
\(\overline y=\Lambda_T^{-1}\sum_{i<T}\lambda_i y_i\) satisfies
\(\|\overline y-z^\star\|^2\le r^2/(2\nu\Lambda_T)\).
Hence the radius is halved after
\begin{equation}\label{eq:base-halving-cost}
T_p(r,\nu)=\max\left\{1,
\left\lceil
\left(\frac{2L_pr^{p-1}}{a_p\nu}\right)^{2/(p+1)}
\right\rceil\right\}
=O_p\!\left(1+\left(\frac{L_pr^{p-1}}{\nu}\right)^{2/(p+1)}\right)
\end{equation}
iterations, each costing at most two queries.

For a supplied radius guess \(\rho>0\), apply these fixed-length
halving stages with radii \(\rho,\rho/2,\ldots\). Stop at the first
radius \(r_J\le(\delta/(A_pL_p))^{1/p}\), then apply
Lemma~\ref{lem:taylor-certificate} and test the evaluated residual.
The predetermined trial cap is
\[
B_0(\rho,\delta)
=2\sum_{j=0}^{J-1}T_p(\rho2^{-j},\nu)+2.
\]
A valid initial radius guarantees success. An invalid radius changes
neither feasibility nor the finite loop lengths; its trial is accepted
only if the evaluated residual is at most \(\delta\).

To avoid requiring the true distance, begin with
\[
\rho_{\min}=\min\{D,(\nu/L_p)^{1/(p-1)}\}
\]
and try \(\rho_{\min},2\rho_{\min},\ldots\), replacing the final
radius by \(D\) and trying \(D\) only once. Each failed trial restarts
from the original supplied point \(w\). If
\(r=\|w-z^\star\|\), the first valid guess \(\rho_*\) satisfies
\[
\rho_*\le\max\{\rho_{\min},2r\},
\qquad
1+\left(\frac{L_p\rho_*^{p-1}}{\nu}\right)^{\gamma_0}
\le C_p\left[1+\left(\frac{L_pr^{p-1}}{\nu}\right)^{\gamma_0}\right].
\]
There are \(O_p(\mathcal L_S)\) guesses and
\(O_p(\mathcal L_S)\) halving stages per guess. The polynomial terms
in the stage costs form a geometric series. Summing all trials up to
\(\rho_*\) proves the claimed bound with logarithmic power two.
The final guess is always valid because \(r\le D\), so the routine is
total. Early exact solutions and successful certificates can only reduce
its cost.
\end{proof}

\subsection{Certified restart as an interface transformation}

A \emph{promised-radius monotone solver} receives a monotone operator
\(G\), a feasible start \(w\), a positive radius \(R\le D\), and a
tolerance \(\delta>0\). Its promise is
\(\operatorname{dist}(w,\operatorname{SOL}(G))\le R\).
It has a finite public query cap and must succeed when the promise is
valid. When the promise is false, it may return failure, but it may report
success only with an evaluated certificate of norm at most \(\delta\).
Suppose its public cap is
\begin{equation}\label{eq:monotone-interface}
C_M\left[1+\left(\frac{L_pR^p}{\delta}\right)^\alpha\right]
\left[1+\ln\!\left(2+\frac DR+\frac{L_pR^p}{\delta}\right)\right]^b.
\end{equation}
Ceilings are understood when an integer cap is installed.

\begin{lemma}[Certified restart]\label{lem:certified-restart}
A promised-radius solver with the bound \eqref{eq:monotone-interface} gives a strong
solver with exponents \((\alpha,b+2)\).
\end{lemma}

\begin{proof}
For a \(\nu\)-strongly monotone operator, the solution \(z^\star\) is
unique. Monotonicity of the normal cone and
\(-G(z^\star)\in N_{\mathcal X}(z^\star)\) imply
\[
\nu\|x-z^\star\|^2
\le\langle G(x)+n,x-z^\star\rangle
\qquad(n\in N_{\mathcal X}(x)).
\]
Consequently,
\begin{equation}\label{eq:restart-error-bound}
\nu\|x-z^\star\|\le\|G(x)+n\|.
\end{equation}

For an initial radius guess \(\rho\), set
\(R_j=\rho2^{-j}\) and
\(\delta_j=\max\{\delta,\nu R_j/2\}\).
Starting at \(w\), invoke the promised-radius solver at stage \(j\)
with radius \(R_j\), tolerance \(\delta_j\), and the public cap in
\eqref{eq:monotone-interface}. If it reports failure, end this radius trial. If it
succeeds and \(\delta_j=\delta\), return its certificate. Otherwise,
use its output as the next stage's start. Equation~\eqref{eq:restart-error-bound}
validates the next radius \(R_{j+1}\). The trial has at most
\[
1+\max\left\{0,
\left\lceil\log_2\frac{\nu\rho}{2\delta}\right\rceil\right\}
\]
stages. A valid initial guess makes every stage valid and guarantees
success. Even after an invalid initial guess, an accepted intermediate
certificate still validates the next radius.

Use the same initial-radius list as in Lemma~\ref{lem:strong-base},
restarting each failed trial from \(w\). Before any original query,
a counter checks the active stage and ancestor budgets; reaching a cap
propagates failure to the radius trial. Thus all failed trials have
predetermined finite costs.
For each stage,
\[
\frac{L_pR_j^p}{\delta_j}
\le\frac{2L_pR_j^{p-1}}{\nu}.
\]
These polynomial terms are geometrically summable. Until termination,
\[
\frac D{R_j}
\le\max\left\{\frac D\rho,\frac{\nu D}{\delta}\right\}
\le\max\left\{1,
\left(\frac{L_pD^{p-1}}{\nu}\right)^{1/(p-1)},
\frac{\nu D}{\delta}\right\}.
\]
It follows that every stage logarithm is at most
\(C_p\mathcal L_S(\nu,\delta)\). One additional logarithm bounds the
number of stages and another bounds the number of radius guesses.
The first valid radius again satisfies
\(\rho_*\le\max\{\rho_{\min},2\|w-z^\star\|\}\).
Summing the stage and trial caps therefore gives the strong-interface
bound with exponents \((\alpha,b+2)\).
\end{proof}

\subsection{A second-difference estimate for Halpern iteration}

The reverse transformation uses the resolvent
\[
P_\eta=(I+\eta(F+N_{\mathcal X}))^{-1}.
\]
For any input \(x\), this is the unique solution of the regularized VI
with operator \(z\mapsto F(z)+(z-x)/\eta\). Existence follows from
compactness and continuity; strong monotonicity gives uniqueness.
If \(u=P_\eta(x)\) and \(v=P_\eta(x')\), monotonicity implies
\[
\langle(x-u)-(x'-v),u-v\rangle\ge0.
\]
Thus \(P_\eta\) is firmly nonexpansive. Its evaluations will be
implemented by the already constructed strong solver.

\begin{lemma}[Halpern second-difference energy]\label{lem:halpern-energy}
Let \(P\) be firmly nonexpansive, with a fixed point \(z^\star\)
satisfying \(\|x_0-z^\star\|\le R\). Define the exact orbit
\[
\overline x_0=x_0,\qquad \overline y_t=P(\overline x_t),\qquad
\overline x_{t+1}=\frac{x_0+(t+1)\overline y_t}{t+2}.
\]
For every integer \(k\ge2\),
\begin{equation}\label{eq:second-difference-energy}
\sum_{t=k}^{2k-1}
\|\overline y_t-2\overline y_{t-1}+\overline y_{t-2}\|^2
\le32\frac{R^2(1+H_{2k})^2}{k^2},
\end{equation}
where \(H_j=\sum_{i=1}^j i^{-1}\). Moreover,
\[
\|\overline x_t-\overline y_t\|
\le\frac{2R(1+H_{t+1})}{t+1}.
\]
\end{lemma}

\begin{proof}
Every \(\overline x_t\) and \(\overline y_t\) lies within \(R\) of
\(z^\star\). Set
\(b_t=\overline x_t-\overline x_{t-1}\) and
\(a_t=\overline y_t-\overline y_{t-1}\).
Subtracting successive updates gives, for \(t\ge2\),
\[
b_t=\frac{t}{t+1}a_{t-1}
       +\frac{\overline y_{t-2}-x_0}{t(t+1)}.
\]
Since \(\|a_t\|\le\|b_t\|\) and \(\|b_1\|\le R\), induction yields
\[
\|b_t\|\le\frac{2RH_t}{t+1}.
\]
Writing \(c_t=b_t-a_{t-1}\), the same recurrence gives
\[
\|c_t\|\le\frac{2R(1+H_{t-1})}{t(t+1)}.
\]
Firm nonexpansiveness gives
\[
\|a_t\|^2+\|b_t-a_t\|^2
\le\|b_t\|^2=\|a_{t-1}+c_t\|^2.
\]
Summing over \(k\le t\le2k-1\) and dropping the final nonnegative
term therefore gives
\[
\sum_{t=k}^{2k-1}\|b_t-a_t\|^2
\le\|a_{k-1}\|^2+
\sum_{t=k}^{2k-1}
       (2\|a_{t-1}\|\|c_t\|+\|c_t\|^2)
\le14\frac{R^2(1+H_{2k})^2}{k^2}.
\]
Also \(\sum_{t=k}^{2k-1}\|c_t\|^2
\le2R^2(1+H_{2k})^2/k^2\).
Since
\(a_t-a_{t-1}=c_t-(b_t-a_t)\), these two estimates prove
\eqref{eq:second-difference-energy}. Finally, for \(t\ge1\),
\[
\overline x_t-\overline y_t
=\frac{x_0-\overline y_{t-1}}{t+1}-a_t.
\]
The displayed residual bound follows from
\(\|x_0-\overline y_{t-1}\|\le2R\) and the bound on \(a_t\).
The case \(t=0\) follows directly from nonexpansiveness.
\end{proof}

\subsection{Extrapolated inexact Halpern iteration and its cost}

\begin{lemma}[Boosting a strong solver]\label{lem:strong-boost}
Suppose a strong solver has exponents \((\gamma,b)\), where
\begin{equation}\label{eq:boost-exponent-range}
\frac{2}{3(p-1)}\le\gamma\le\frac{2}{p+1}.
\end{equation}
Then there is a promised-radius monotone solver of the form
\eqref{eq:monotone-interface} with exponents
\(\alpha=\gamma/(1+\gamma)\) and \(b+4\).
\end{lemma}

\begin{proof}
Consider a monotone problem with start \(x_0\), positive radius promise
\(R\le D\), and target residual \(\varepsilon\). Define
\[
Q_R=\frac{L_pR^p}{\varepsilon},\qquad
\mathcal L_A=1+\ln\!\left(2+\frac DR+Q_R\right),\qquad
K=\max\{2,\lceil Q_R^\alpha\rceil\}.
\]
Use the parameters
\begin{equation}\label{eq:halpern-parameters}
e=\frac{R}{(K+2)^4},\qquad
\eta=\frac{8R(1+H_{K+1})}{\varepsilon(K+1)},\qquad
\delta=\frac e\eta.
\end{equation}
At iteration \(t\), define
\(G_t(z)=F(z)+(z-x_t)/\eta\), which is \(1/\eta\)-strongly monotone
and retains the original \(L_p\)-smoothness promise. Its exact solution
is \(y_t^*=P_\eta(x_t)\).

The implemented algorithm has \(K+1\) inner calls:
\begin{enumerate}[label=\arabic*.,leftmargin=*]
\item Set \(x_0\) equal to the supplied start. At \(t=0\), use
\(w_0=x_0\); at \(t=1\), use \(w_1=\widehat y_0\); and, for
\(t\ge2\), use the feasible extrapolation
\[
w_t=\Pi_{\mathcal X}(2\widehat y_{t-1}-\widehat y_{t-2}).
\]
\item Invoke the input strong solver on \(G_t\), starting from \(w_t\),
with modulus \(1/\eta\) and tolerance \(\delta\). Save its evaluated
certificate \((\widehat y_t,n_t,G_t(\widehat y_t))\).
\item If \(t<K\), set
\[
x_{t+1}=\frac{x_0+(t+1)\widehat y_t}{t+2}
\]
and continue. After \(t=K\), recover
\(F(\widehat y_K)=G_K(\widehat y_K)-(\widehat y_K-x_K)/\eta\).
Return the saved certificate only if its evaluated norm
\(\|F(\widehat y_K)+n_K\|\) is at most \(\varepsilon\); otherwise
report failure.
\end{enumerate}
A shared raw-query cap, specified below, aborts a run before its next
query would exceed the budget. An affine-regularized jet costs exactly
one original jet. Projections and convex combinations keep every query
feasible.

Assume first that the radius promise is valid, and choose a solution
\(z^\star\) with \(\|x_0-z^\star\|\le R\). Let
\((\overline x_t,\overline y_t)\) be the exact shadow orbit of
Lemma~\ref{lem:halpern-energy} for \(P_\eta\), with the same anchor.
Strong monotonicity and the inner certificate imply
\[
\|\widehat y_t-y_t^*\|\le\eta\delta=e.
\]
Induction through the two Halpern recursions gives
\[
\|x_t-\overline x_t\|\le\frac{te}{2},\qquad
\|\widehat y_t-\overline y_t\|\le\frac{(t+2)e}{2}.
\]
For the actual warm-start distances \(r_t=\|w_t-y_t^*\|\),
nonexpansiveness of projection and of \(P_\eta\) gives
\(r_0,r_1\le3R\), and, for \(t\ge2\),
\[
r_t\le
\|\overline y_t-2\overline y_{t-1}+\overline y_{t-2}\|
+(2t+1)e.
\]

Put \(a=(p-1)\gamma\); condition~\eqref{eq:boost-exponent-range} implies
\(2/3\le a<2\). Power mean and
Lemma~\ref{lem:halpern-energy} give, on a dyadic block of length \(k\),
\[
\sum_{t=k}^{2k-1}
\|\overline y_t-2\overline y_{t-1}+\overline y_{t-2}\|^a
\le32^{a/2}R^a(1+H_{2k})^a k^{1-3a/2}.
\]
The exponent of \(k\) is nonpositive. Summing the dyadic blocks and
using \((u+v)^a\le2(u^a+v^a)\) yields
\begin{equation}\label{eq:warm-start-sum}
\sum_{t=0}^{K}r_t^{(p-1)\gamma}
\le C_pR^{(p-1)\gamma}
            [1+\ln(K+2)]^{(p-1)\gamma+1}.
\end{equation}
Indeed, the accumulated perturbation contributes at most
\[
2e^a\sum_{t=2}^{K}(2t+1)^a
\le2\cdot3^aR^a K^{1-3a}\le2\cdot3^aR^a.
\]
The exact shadow orbit may be extended beyond \(K\) when bounding an
incomplete final block; no additional queries are performed.

For correctness of the returned point, the saved regularized residual
and the preceding orbit estimates imply
\begin{equation}\label{eq:boost-certificate}
\|F(\widehat y_K)+n_K\|
\le\eta^{-1}\left[
\frac{2R(1+H_{K+1})}{K+1}+(K+2)e\right]
<\varepsilon.
\end{equation}
The first term after division by \(\eta\) is \(\varepsilon/4\),
and the second is smaller than \(\varepsilon/8\). Thus every
valid-promise run passes the final evaluated test.

We next count all inner calls. Let \(h=1+H_{K+1}\).
The logarithmic factor for every inner strong solve is
\(\mathcal L_S(1/\eta,e/\eta)\). Its three dimensionless terms are
\[
\eta L_pD^{p-1}
=\frac{8hQ_R}{K+1}\left(\frac DR\right)^{p-1},
\]
\[
\frac{L_pD^p}{e/\eta}
=\frac{8hQ_R(K+2)^4}{K+1}\left(\frac DR\right)^p,
\qquad
\frac{(1/\eta)D}{e/\eta}=\frac DR(K+2)^4.
\]
Since \(K\) is polynomially bounded by \(1+Q_R\), these identities
show that
\(\mathcal L_S(1/\eta,e/\eta)\le C_p\mathcal L_A\), while
\(h\) and \(1+\ln(K+2)\) are also at most \(C_p\mathcal L_A\).
Summing the input strong-interface costs using
\eqref{eq:warm-start-sum} gives a bound with logarithmic power
\(b+1+p\gamma\). By \eqref{eq:boost-exponent-range},
\(1+p\gamma\le1+2p/(p+1)<3\); the following slightly larger power
is convenient for the whole hierarchy:
\begin{equation}\label{eq:boost-query-cost}
C_p\mathcal L_A^{\,b+4}
\left[K+1+\left(\frac{Q_R}{K+1}\right)^\gamma\right].
\end{equation}
The identity \(\gamma(1-\alpha)=\alpha\) and the definition of
\(K\) show that its bracket is \(O_p(1+Q_R^\alpha)\).

For completeness, the following entirely public finite sum implements
the cap used in this argument:
\[
U_a(K)=4\cdot3^a+
2\cdot32^{a/2}\sum_{j=1}^{\lfloor\log_2K\rfloor}
 (1+H_{2^{j+1}})^a\,2^{j(1-3a/2)}.
\]
The preceding block calculation gives
\(\sum_{t=0}^{K}r_t^a\le R^aU_a(K)\).
If \(C_S\) is the fixed constant of the input strong solver, install
the shared original-oracle cap
\[
B=\left\lceil
C_S\mathcal L_S(1/\eta,e/\eta)^b
\left[K+1+(\eta L_p)^\gamma R^aU_a(K)\right]
\right\rceil+1.
\]
On a valid promise, the sum of all completed inner-solver costs is at
most \(B\), so the cap cannot interrupt the run. On an invalid
promise, every descendant checks this same counter before making a
query. An exhausted budget propagates failure; a completed run must
still pass the final evaluated residual test. This proves the claimed
promised-radius interface, including finite caps and sound acceptance.
\end{proof}

\subsection{Finite acceleration and completion of the proof}

Start with Lemma~\ref{lem:strong-base}, then apply
Lemma~\ref{lem:strong-boost} followed by
Lemma~\ref{lem:certified-restart}. At hierarchy level \(j\), the strong
solver has exponents
\begin{equation}\label{eq:finite-hierarchy}
\gamma_j=\frac{2}{p+1+2j},
\qquad b_j=6j+2.
\end{equation}
This follows by induction from
\(\gamma\mapsto\gamma/(1+\gamma)\), with four logarithmic powers
assigned to the boost and two to restart. For each
\(j=1,\ldots,p-1\), the input exponent satisfies
\[
\frac{2}{3(p-1)}\le\gamma_{j-1}\le\frac{2}{p+1},
\qquad
(p-1)\gamma_{j-1}
=\frac{2(p-1)}{p+2j-1}\ge\frac23.
\]
Every invocation of the boost is therefore valid, including the endpoint
in the last layer. At \(j=p-1\), equation~\eqref{eq:finite-hierarchy} gives
\(\gamma_{p-1}=2/(3p-1)\) and \(b_{p-1}=6p-4\).
The hierarchy has exactly \(p-1\) layers; each layer uses only routines
already constructed at lower levels.

Apply the final strong solver to \(F\), with start \(x_0\), modulus
\(\mu\), and tolerance \(\varepsilon\). Its distance-sensitive bound
is \eqref{eq:distance-sensitive-upper}. Since \(r_0=\|x_0-x^\star\|\le D\) and
\begin{equation}\label{eq:log-identities}
\frac{L_pD^p}{\varepsilon}=\frac{\kappa_D}{\theta},
\qquad
\frac{\mu D}{\varepsilon}=\frac1\theta,
\end{equation}
we have
\(\mathcal L_S(\mu,\varepsilon)=\mathcal H(\kappa_D,\theta)\).
This proves \eqref{eq:hierarchy-upper}.

Finally, the construction respects the policy requirements of
Section~\ref{sec:setup}. For each fixed public domain, known-coefficient
model VIs have unique solutions on every legal transcript. Compactness
and passage to the VI limit show that these solutions depend continuously
on the coefficients within the unique-solution class. The closed,
compact-valued solution relation admits a Borel selection on the remaining
coefficient tuples, which makes the rules total there as well.
Tensor arithmetic, projections, finite loops, and evaluated norm tests
are measurable. Radius lists, halving stages, and ancestor counters are
finite and public. Every accepted certificate uses an operator value at
the returned point, and affine corrections require no additional query.
These observations complete the proof of Theorem~\ref{thm:upper}.

\section{Randomized Lower Bound}\label{sec:lower}

The hard instance is a smooth contractive chain placed in an unknown
orthonormal frame. Its flat gates hide complete derivative tensors, while
the contraction defect determines the strong-monotonicity modulus. We
first establish the geometric and analytic properties of the chain, then
prove that arbitrary adaptive queries cannot uncover it quickly enough.

\subsection{A smooth chain with flat gates}

Let
\[
\eta_0(t)=
\begin{cases}
e^{-1/t},&t>0,\\
0,&t\le0,
\end{cases}
\qquad
\chi(t)=\frac{\eta_0(t)}{\eta_0(t)+\eta_0(1-t)},
\qquad
\varrho(u)=\chi\!\left(\frac{4u^2-1}{3}\right).
\]
Thus $\varrho$ is an even $C^\infty$ function with values in $[0,1]$,
equal to zero on $[-1/2,1/2]$ and to one outside $(-1,1)$. Define
\begin{equation}\label{eq:flat-gate}
B_p=\max\{1,\|\varrho^{(p-1)}\|_\infty\},
\qquad
h_\tau(s)=\tau\int_0^{s/\tau}\varrho(u)\,du.
\end{equation}

\begin{lemma}[Flat-gate properties]\label{lem:flat-gate}
For every $\tau>0$, the function $h_\tau$ is odd and $C^\infty$.
It and all its derivatives vanish on $[-\tau/2,\tau/2]$. Moreover,
\[
0\le h_\tau'(s)\le1,
\qquad |h_\tau(s)|\le |s|,
\qquad |s-h_\tau(s)|\le\tau,
\]
and
\begin{equation}\label{eq:gate-smoothness}
\operatorname{Lip}(h_\tau^{(p-1)})
\le B_p\tau^{1-p}.
\end{equation}
\end{lemma}

\begin{proof}
Oddness follows from the evenness of $\varrho$, and
$h_\tau'(s)=\varrho(s/\tau)$. Exact flatness follows from the
definition on the open slab and smoothness at its endpoints.
Since $0\le\varrho\le1$, the function $h_\tau$ is 1-Lipschitz and
vanishes at zero. For $s\ge0$,
\[
0\le s-h_\tau(s)
=\int_0^s(1-\varrho(t/\tau))\,dt\le\tau;
\]
oddness gives the same absolute bound for $s<0$. Finally,
$h_\tau^{(p)}(s)=\tau^{1-p}\varrho^{(p-1)}(s/\tau)$, which proves
\eqref{eq:gate-smoothness}.
\end{proof}

For an integer $m\ge2$, set
\begin{equation}\label{eq:chain-parameters}
R=\frac D2,\qquad
q=1-\frac1{2m},\qquad
c=\frac{R}{2\sqrt m},\qquad
\tau=\frac{R}{16m^{3/2}},\qquad
\lambda=\frac{L_p\tau^{p-1}}{B_p}.
\end{equation}
Let $\mathcal X=\{x\in\R^d:\|x\|\le R\}$ and $x_0=0$.
For an orthonormal $m$-frame $V=(v_1,\ldots,v_m)$, define
\begin{equation}\label{eq:hard-family}
T_V(x)=cv_1+q\sum_{i=1}^{m-1}
h_\tau(\langle v_i,x\rangle)v_{i+1},
\qquad
F_V(x)=\lambda(x-T_V(x)).
\end{equation}

\begin{lemma}[Geometry and regularity of the chain]\label{lem:chain-regularity}
The map $T_V$ is a contraction and preserves $\mathcal X$:
\begin{equation}\label{eq:chain-contraction}
\|T_V(x)-T_V(y)\|\le q\|x-y\|,
\qquad T_V(\mathcal X)\subseteq\mathcal X.
\end{equation}
The operator $F_V$ is globally smooth and strongly monotone with
certified modulus
\begin{equation}\label{eq:chain-modulus}
\mu_m=\lambda(1-q)=\frac{\lambda}{2m}.
\end{equation}
Its highest-derivative Lipschitz modulus is at most $L_p$.
Its unique zero $x^\star$ is interior to $\mathcal X$ and satisfies
\begin{equation}\label{eq:interior-zero}
\frac{3R}{16}\le\|x^\star\|\le\frac R2<R.
\end{equation}
\end{lemma}

\begin{proof}
By orthogonality, the 1-Lipschitz property of $h_\tau$, and Bessel's
inequality,
\[
\|T_V(x)-T_V(y)\|^2
\le q^2\sum_{i=1}^{m-1}|\langle v_i,x-y\rangle|^2
\le q^2\|x-y\|^2.
\]
For $\|x\|\le R$, the same calculation gives
\[
\|T_V(x)\|^2\le c^2+q^2R^2
=R^2\left(1-\frac{3m-1}{4m^2}\right)\le R^2.
\]
The contraction bound implies
\[
\langle F_V(x)-F_V(y),x-y\rangle
\ge\lambda(1-q)\|x-y\|^2,
\]
which proves \eqref{eq:chain-modulus}.

To verify smoothness without a dimension factor, take arbitrary unit
vectors $u_1,\ldots,u_p$. Since $p\ge2$, the affine part of $F_V$
has zero $p$th derivative. The nonlinear part satisfies
\begin{align*}
\|D^pF_V(x)[u_1,\ldots,u_p]\|^2
&=\lambda^2q^2\sum_{i=1}^{m-1}
 |h_\tau^{(p)}(\langle v_i,x\rangle)|^2
 \prod_{j=1}^p|\langle v_i,u_j\rangle|^2\\
&\le\lambda^2q^2B_p^2\tau^{2-2p}
 \sum_{i=1}^{m-1}|\langle v_i,u_1\rangle|^2
\le\lambda^2q^2B_p^2\tau^{2-2p}.
\end{align*}
Consequently,
\begin{equation}\label{eq:chain-smoothness}
\|D^pF_V(x)\|_{\mathrm{op}}
\le\lambda qB_p\tau^{1-p}\le L_p.
\end{equation}
Integration along line segments gives the required Lipschitz bound on
$D^{p-1}F_V$.

The contraction has a unique fixed point
$x^\star=\sum_{i=1}^m s_i v_i$, with
$s_1=c$ and $s_{i+1}=qh_\tau(s_i)$. Since
$0\le h_\tau(s)\le s$ for $s\ge0$,
$0\le s_i\le cq^{i-1}\le c$. The estimate
$h_\tau(s)\ge s-\tau$ also gives
\[
s_i\ge cq^{i-1}-\tau\sum_{j=1}^{i-1}q^j
\ge\frac c2-m\tau=\frac{3c}{8}.
\]
Here $q^{m-1}\ge1/2$ by Bernoulli's inequality and $m\tau=c/8$.
Thus $3c\sqrt m/8\le\|x^\star\|\le c\sqrt m$, proving
\eqref{eq:interior-zero}.
\end{proof}

\subsection{A boundary-safe residual obstruction}

\begin{lemma}[A hidden final coordinate forces a large residual]
\label{lem:hidden-residual}
For the operator \eqref{eq:hard-family},
\begin{equation}\label{eq:boundary-residual}
r_{\mathrm{tan}}(x)=\|F_V(x)\|
\qquad(x\in\mathcal X).
\end{equation}
If $|\langle v_m,x\rangle|\le\tau/2$, then
\begin{equation}\label{eq:hidden-slab-residual}
r_{\mathrm{tan}}(x)
\ge\frac{\lambda R}{8m}
=\frac{\mu_mD}{8}.
\end{equation}
\end{lemma}

\begin{proof}
The normal cone is zero in the interior. On the boundary,
$N_{\mathcal X}(x)=\{ax:a\ge0\}$ and, by \eqref{eq:chain-contraction},
\[
\langle F_V(x),x\rangle
=\lambda(R^2-\langle T_V(x),x\rangle)\ge0.
\]
It follows that $\|F_V(x)+ax\|^2\ge\|F_V(x)\|^2$ for all $a\ge0$,
so \eqref{eq:boundary-residual} also holds at the boundary.

Write $s_i=\langle v_i,x\rangle$, $e=x-T_V(x)$, and
$w=\sum_{i=1}^m q^{m-i}v_i$. The chain telescopes exactly:
\begin{align*}
\langle w,e\rangle
&=s_m-cq^{m-1}
 +\sum_{i=1}^{m-1}q^{m-i}\bigl(s_i-h_\tau(s_i)\bigr)\\
&\le\frac\tau2-\frac c2+m\tau
\le-\frac c4.
\end{align*}
We used $|s_m|\le\tau/2$, Lemma~\ref{lem:flat-gate},
$m\tau=c/8$, and $\tau/2\le c/16$.
Since $\|w\|\le\sqrt m$, Cauchy--Schwarz gives
$\|e\|\ge c/(4\sqrt m)=R/(8m)$. Multiplying by $\lambda$ and using
\eqref{eq:chain-modulus} and \eqref{eq:boundary-residual} proves \eqref{eq:hidden-slab-residual}.
\end{proof}

\subsection{Exact hiding under adaptive randomized queries}

The next lemma is the information-theoretic step. It concerns one
distribution of fixed instances, drawn independently of the algorithm's
seed, and includes the algorithm's final output in the hiding event.

\begin{lemma}[Adaptive full-jet hiding]\label{lem:adaptive-hiding}
Let $N=m-1$ and suppose
\begin{equation}\label{eq:lower-dimension}
d\ge m+
\left\lceil2048m^3\ln\bigl(8(m^2+1)\bigr)\right\rceil.
\end{equation}
For any randomized adaptive policy making at most $N$ feasible full-jet
queries and returning a feasible point, there is a fixed frame $V_0$
such that
\[
\Pr_{\mathrm{seed}}\!\left\{
r_{\mathrm{tan}}^{F_{V_0}}(\widehat x)
\ge\frac{\mu_mD}{8}\right\}\ge\frac34.
\]
The frame may depend on the policy, but is independent of its realized
seed.
\end{lemma}

\begin{proof}
Pad early stopping with ignored queries while retaining the saved output.
We use total Borel query and output rules on finite transcripts. Applying
the known projection onto $\mathcal X$ and enforcing the cap $N$ extends
feasibility and stopping to all transcripts without changing any legal
run. This makes the following auxiliary simulation well defined even
when its replies do not form a legal transcript.

Draw the complete frame $V$ from Haar measure, independently of the seed
$U$. At simulated query $t\in\{1,\ldots,N\}$, return the full jet of
\[
F_t^V(x)=\lambda\left(
x-cv_1-q\sum_{i=1}^{t-1}
h_\tau(\langle v_i,x\rangle)v_{i+1}\right).
\]
Let $\widetilde x_t$ and $\widetilde x_{\mathrm{out}}$ be the queries
and saved output in this simulation. Inductively,
$\widetilde x_t$ is measurable with respect to
$(U,v_1,\ldots,v_{t-1})$, and the response at step $t$ depends only on
$(U,v_1,\ldots,v_t)$. The final output is therefore measurable with
respect to $(U,v_1,\ldots,v_N)$. This argument does not require the
policy to recover the frame vectors from its observations.

We use the following elementary spherical tail bound. If $v$ is uniform
on the unit sphere of a $k$-dimensional subspace and $\|x\|\le R$, then
\begin{equation}\label{eq:spherical-tail}
\Pr\{|\langle v,x\rangle|>a\}
\le2\exp\!\left(-\frac{ka^2}{2R^2}\right).
\end{equation}
Indeed, with $x$ first projected onto that subspace, the even moments
satisfy
\[
\mathbb E\langle v,x\rangle^{2j}
=\frac{(2j-1)!!\,\|x\|^{2j}}
 {k(k+2)\cdots(k+2j-2)}
\le\frac{(2j-1)!!\,R^{2j}}{k^j}.
\]
Symmetry and the exponential series imply
$\mathbb E e^{t\langle v,x\rangle}\le e^{t^2R^2/(2k)}$.
Applying the exponential Markov inequality to both signs proves
\eqref{eq:spherical-tail}.

Define the good event
\[
\mathcal E=
\bigcap_{t=1}^{N}\ \bigcap_{i=t}^{m-1}
\left\{|\langle v_i,\widetilde x_t\rangle|\le\frac\tau2\right\}
\ \cap\
\left\{|\langle v_m,\widetilde x_{\mathrm{out}}\rangle|
\le\frac\tau2\right\}.
\]
Conditioning on $(U,v_1,\ldots,v_{t-1})$ fixes
$\widetilde x_t$. Each remaining $v_i$ has a uniform spherical marginal
in the orthogonal complement of this prefix, whose dimension is
$d-t+1$. For the final-output event, condition instead on
$(U,v_1,\ldots,v_N)$; the remaining dimension is $d-m+1$.
In each case the dimension is at least $d-m$. We apply
\eqref{eq:spherical-tail} under this conditioning and then integrate, without
conditioning on any earlier slab event. There are
$m(m-1)/2+1<m^2+1$ events. Since $\tau/2=R/(32m^{3/2})$, the union
bound and \eqref{eq:lower-dimension} yield
\[
\Pr(\mathcal E^c)
\le2(m^2+1)\exp\!\left(-\frac{d-m}{2048m^3}\right)
\le\frac14.
\]

On $\mathcal E$, every gate omitted from $F_t^V$ and every derivative
of that gate vanishes at $\widetilde x_t$. Thus
$\mathcal O_p^{F_V}(\widetilde x_t)
=\mathcal O_p^{F_t^V}(\widetilde x_t)$. Induction on $t$ shows that
the genuine and simulated transcripts, queries, and saved outputs agree
exactly. Lemma~\ref{lem:hidden-residual} then implies, jointly over the
frame and seed,
\begin{equation}\label{eq:joint-failure}
\Pr_{V,U}\!\left\{
r_{\mathrm{tan}}^{F_V}(\widehat x)
\ge\frac{\lambda R}{8m}\right\}\ge\frac34.
\end{equation}
All events are measurable: the hard-family jets are continuous in the
frame and query point, and the finite policy recursion is Borel.
Averaging \eqref{eq:joint-failure} over frames gives a fixed $V_0$ for which
the same failure probability over $U$ is at least $3/4$.
\end{proof}

\subsection{Calibration to the public monotonicity modulus}

We now prove Theorem~\ref{thm:lower}. Define constants depending only on
$p$ by
\begin{equation}\label{eq:lower-constants}
c_p^{\mathrm{sc}}=\frac1{2^pB_p16^{p-1}},\qquad
\kappa_{0,p}=\frac{2^{(3p-1)/2}}{c_p^{\mathrm{sc}}},\qquad
c_p=\frac12(c_p^{\mathrm{sc}})^{2/(3p-1)}.
\end{equation}
For any public tuple satisfying \eqref{eq:lower-regime}, choose
\begin{equation}\label{eq:chain-length}
s=(c_p^{\mathrm{sc}}\kappa_D)^{2/(3p-1)},
\qquad m=\lfloor s\rfloor.
\end{equation}
Then $s\ge2$, so $m\ge2$ and $m\ge s/2$. Substituting
\eqref{eq:chain-parameters} into \eqref{eq:chain-modulus} gives
\begin{equation}\label{eq:modulus-calibration}
\mu_m
=c_p^{\mathrm{sc}}\frac{L_pD^{p-1}}{m^{(3p-1)/2}}
=\mu\left(\frac{s}{m}\right)^{(3p-1)/2}
\ge\mu.
\end{equation}
Together with Lemma~\ref{lem:chain-regularity}, this verifies every
public promise for every frame: the domain has diameter $D$, the
highest-derivative Lipschitz modulus is at most $L_p$, and the
strong-monotonicity modulus is at least $\mu$.

Choose a dimension satisfying \eqref{eq:lower-dimension}. If a randomized policy
has cap $m-1$, Lemma~\ref{lem:adaptive-hiding} provides one fixed legal
instance on which, with probability at least $3/4$,
\begin{equation}\label{eq:fixed-instance-failure}
r_{\mathrm{tan}}(\widehat x)
\ge\frac{\mu_mD}{8}
\ge\frac{\mu D}{8}
\ge2\varepsilon.
\end{equation}
Its success probability is therefore at most $1/4$, contradicting the
required $2/3$. Hence
\[
\Trand(\varepsilon;L_p,\mu,D)
\ge m\ge\frac s2
=c_p\kappa_D^{2/(3p-1)},
\]
as claimed. The polynomial exponent comes directly from the chain's
calibration
$\mu_m\asymp_p L_pD^{p-1}m^{-(3p-1)/2}$; the random frame removes
restrictions on the adaptive policy without changing this exponent.


\section{Transfer to Scalar Saddle Problems}\label{sec:saddle}

The upper bound also applies to strongly convex--strongly concave
objectives. The transfer preserves the smoothness constant, the
residual, and the number of oracle calls.

Let $\mathcal X\subseteq\R^{d_x}$ and
$\mathcal Y\subseteq\R^{d_y}$ be known nonempty compact convex sets,
and suppose $\mathcal Z=\mathcal X\times\mathcal Y$ has diameter at
most $D$ in the Euclidean product norm
$\|(x,y)\|^2=\|x\|^2+\|y\|^2$. A feasible start $z_0$ is public.
Let $\phi$ be $C^p$ on a neighborhood of $\mathcal Z$, with
$\mu_x$-strongly convex $x$-sections and $\mu_y$-strongly concave
$y$-sections, where $\mu_x,\mu_y>0$. Assume
\begin{equation}\label{eq:scalar-smoothness}
\|D^p\phi(z)-D^p\phi(z')\|_{\mathrm{sc},p}
\le L_p\|z-z'\|
\qquad(z,z'\in\mathcal Z),
\end{equation}
where $\|\cdot\|_{\mathrm{sc},p}$ is the norm of a scalar $p$-linear
form, taking the supremum of its absolute value over unit inputs.
One scalar oracle call returns
\begin{equation}\label{eq:scalar-jet}
(\phi(z),D\phi(z),\ldots,D^p\phi(z)).
\end{equation}
All parameter promises and policy conventions are as in
Section~\ref{sec:setup}.

Write $\mathsf T_{p,\mathrm{sad}}^{\mathrm{det,cert}}$ for the
corresponding deterministic certificate complexity. The required
output is an evaluated triple $(\widehat z,n,\nabla\phi(\widehat z))$,
where $n=(n_x,n_y)\in
N_{\mathcal X}(\widehat x)\times N_{\mathcal Y}(\widehat y)$ and
\begin{equation}\label{eq:saddle-certificate}
\left\|
\binom{\nabla_x\phi(\widehat z)}{-\nabla_y\phi(\widehat z)}+n
\right\|\le\varepsilon.
\end{equation}
The gradient in this certificate must come from a charged call at
$\widehat z$.

\begin{theorem}[Exact scalar-jet simulation]\label{thm:transfer}
With $\mu=\min\{\mu_x,\mu_y\}$,
\begin{equation}\label{eq:scalar-simulation}
\mathsf T_{p,\mathrm{sad}}^{\mathrm{det,cert}}
(\varepsilon;L_p,\mu_x,\mu_y,D)
\le
\Tcert(\varepsilon;L_p,\mu,D).
\end{equation}
Each operator call is simulated by exactly one scalar call. The
simulation preserves all nested and final evaluations, and the
operator tangent residual equals the product KKT residual.
\end{theorem}

\begin{proof}
On $E=\R^{d_x}\times\R^{d_y}$, set
\begin{equation}\label{eq:signed-gradient}
J=\operatorname{diag}(I_{d_x},-I_{d_y}),
\qquad
F_\phi=J\nabla\phi
=\binom{\nabla_x\phi}{-\nabla_y\phi}.
\end{equation}
Add the strong-convexity inequalities in $x$ at
$(x_1,y_1)$ and $(x_2,y_2)$, evaluated respectively at $x_2$ and
$x_1$, and the strong-concavity inequalities at the same two points,
evaluated respectively at $y_2$ and $y_1$. Their four function values
cancel, giving
\begin{equation}\label{eq:saddle-monotonicity}
\langle F_\phi(z_1)-F_\phi(z_2),z_1-z_2\rangle
\ge\mu_x\|x_1-x_2\|^2+\mu_y\|y_1-y_2\|^2
\ge\mu\|z_1-z_2\|^2.
\end{equation}
This proof uses only feasible chords, so it includes boundary points
and lower-dimensional factors.

For a scalar $(k+1)$-linear form $A$, define the vector-valued
$k$-linear form $\mathcal R_k(A)$ by Euclidean duality:
\[
\langle\mathcal R_k(A)[u_1,\ldots,u_k],u_0\rangle
=A[u_0,u_1,\ldots,u_k].
\]
For $k=0$, this is the usual identification of a linear functional
with its gradient vector. Since $J$ is orthogonal,
\begin{equation}\label{eq:tensor-isometry}
\|J\mathcal R_k(A)\|_{\mathrm{vec},k}
=\|A\|_{\mathrm{sc},k+1}.
\end{equation}
Furthermore, for $0\le k\le p-1$,
\begin{equation}\label{eq:jet-conversion}
D^kF_\phi(z)=J\mathcal R_k(D^{k+1}\phi(z)).
\end{equation}
Equations \eqref{eq:scalar-smoothness}, \eqref{eq:tensor-isometry}, and \eqref{eq:jet-conversion}
give $\operatorname{Lip}(D^{p-1}F_\phi)\le L_p$ with no dimension
factor. Thus $F_\phi$ belongs to the operator class with public
parameters $(p,L_p,\mu,D)$.

The outward normal cone factors as
\begin{equation}\label{eq:product-normal}
N_{\mathcal Z}(x,y)
=N_{\mathcal X}(x)\times N_{\mathcal Y}(y).
\end{equation}
Consequently, the constrained first-order saddle conditions are
equivalent to $0\in F_\phi(z)+N_{\mathcal Z}(z)$, and
\begin{equation}\label{eq:residual-identity}
r_{\mathrm{KKT}}(x,y)
:=\operatorname{dist}\!\left(0,
\binom{\nabla_x\phi(x,y)}{-\nabla_y\phi(x,y)}
+N_{\mathcal X}(x)\times N_{\mathcal Y}(y)\right)
=r_{\mathrm{tan}}^{F_\phi}(x,y).
\end{equation}

To implement any deterministic operator policy, make one scalar call
at each point it requests and transform the returned tensors using
\eqref{eq:jet-conversion}. These are fixed linear transformations. The converted
responses are exactly the operator's full jets, so the two transcripts
and all subsequent queries agree. This includes every nested call,
the stopping rule, and the final certificate evaluation. Equations
\eqref{eq:product-normal}--\eqref{eq:residual-identity} convert the saved certificate without
an additional query. The policy therefore has the same pathwise cap,
which proves \eqref{eq:scalar-simulation}.
\end{proof}

Combining Theorems~\ref{thm:upper} and~\ref{thm:transfer} yields
\begin{equation}\label{eq:saddle-hierarchy-upper}
\mathsf T_{p,\mathrm{sad}}^{\mathrm{det,cert}}
(\varepsilon;L_p,\mu_x,\mu_y,D)
\le C_p(1+\kappa_D^{2/(3p-1)})
\mathcal H(\kappa_D,\theta)^{6p-4},
\qquad \mu=\min\{\mu_x,\mu_y\}.
\end{equation}
The sharper consequence of Proposition~\ref{prop:literature-restart}
also transfers without loss:
\begin{equation}\label{eq:saddle-sharper}
\mathsf T_{p,\mathrm{sad}}^{\mathrm{det,cert}}
(\varepsilon;L_p,\mu_x,\mu_y,D)
\le C_p\left(1+\kappa_D^{\ap}+\logplus\frac1\theta\right).
\end{equation}
As in that proposition, this estimate explicitly uses the monotone
solver of~\cite{zhangchen2026}.

\begin{corollary}[From a KKT certificate to a saddle gap]
\label{cor:saddle-gap}
For every $z=(x,y)\in\mathcal Z$, define
\[
\operatorname{Gap}_\phi(x,y)
=\max_{y'\in\mathcal Y}\phi(x,y')
-\min_{x'\in\mathcal X}\phi(x',y).
\]
Then, with $\mu=\min\{\mu_x,\mu_y\}$,
\begin{equation}\label{eq:saddle-gap}
\operatorname{Gap}_\phi(z)
\le\min\!\left\{D\,r_{\mathrm{KKT}}(z),
\frac{r_{\mathrm{KKT}}(z)^2}{2\mu}\right\}.
\end{equation}
\end{corollary}

\begin{proof}
Fix $z'=(x',y')\in\mathcal Z$ and $n\in N_{\mathcal Z}(z)$.
Strong convexity and concavity, followed by the normal-cone inequality,
give
\begin{align*}
\phi(x,y')-\phi(x',y)
&\le\langle F_\phi(z),z-z'\rangle
 -\frac\mu2\|z-z'\|^2\\
&\le\langle F_\phi(z)+n,z-z'\rangle
 -\frac\mu2\|z-z'\|^2.
\end{align*}
The last expression is at most $D\|F_\phi(z)+n\|$ by the diameter
bound, and at most $\|F_\phi(z)+n\|^2/(2\mu)$ by completing the
square. Take the supremum over $z'$ and then the infimum over $n$.
\end{proof}

The transfer in this section is an upper-bound result. A matching
lower bound for scalar saddle problems would require a hard family
with the signed-gradient structure in \eqref{eq:signed-gradient}; the general
operator construction in Section~\ref{sec:lower} does not supply that
property. Likewise, \eqref{eq:saddle-gap} is a one-way implication
and does not by itself yield a gap lower bound.


\section{Discussion}\label{sec:discussion}
The results identify the condition dependence of exact higher-order
information for strongly monotone VIs. The lower bound applies to
arbitrary adaptive randomized policies, whereas both upper bounds return
the stronger evaluated certificate. At fixed relative accuracy, the
self-contained hierarchy matches the polynomial condition exponent;
the restart consequence of~\cite{zhangchen2026} removes the remaining
logarithmic condition loss. This distinction records both the complete
proof given here and the strongest comparison supported by the current
literature.

\paragraph{Accuracy and small condition parameters.}
The upper bounds hold for every $\kappa_D>0$ and $\theta>0$.
The lower bound requires $\kappa_D\ge\kappa_{0,p}$ and
$\theta\le1/16$. It does not establish a matching logarithmic accuracy
term, and thus does not characterize the joint dependence on
$(\kappa_D,\theta)$ in every regime. In particular, neither the factor
$\mathcal H^{6p-4}$ in the self-contained proof nor the term
$\logplus(1/\theta)$ in the imported restart bound is claimed to be
necessary. Sharper local convergence and an accuracy-sensitive lower
bound are separate questions.

\paragraph{Dimension and computation.}
The upper query guarantees are uniform over finite dimensions. The
randomized lower bound uses a dimension of order $m^3\log m$, with
$m\asymp_p\kappa_D^{\ap}$; it makes no fixed-dimensional optimality
claim. Its concentration argument also uses the bounded feasible-query
domain. Off-domain, noisy, or partial-tensor oracles require a separate
analysis. Finally, an exact full tensor is charged as one reply, and
Taylor-model VIs formed from known coefficients are solved by internal
computation. Converting these information bounds into arithmetic or bit
complexity requires additional assumptions on the domain representation
and the model solver.

\paragraph{Scalar structure.}
The saddle reduction gives the same upper guarantees with
$\mu=\min\{\mu_x,\mu_y\}$, and a KKT certificate controls both distance
and saddle gap. The randomized lower construction is a general
strongly monotone operator; it does not have the signed-gradient
structure needed for a scalar saddle lower bound. Identifying the
optimal accuracy dependence, the smallest sufficient dimension, and a
matching unrestricted scalar lower family are natural extensions of
the condition-complexity statement proved here.

\begingroup
\small
\setlength{\bibsep}{2pt plus 1pt minus 1pt}
\urlstyle{same}
\begin{thebibliography}{19}
\providecommand{\natexlab}[1]{#1}
\providecommand{\url}[1]{\texttt{#1}}
\expandafter\ifx\csname urlstyle\endcsname\relax
  \providecommand{\doi}[1]{doi: #1}\else
  \providecommand{\doi}{doi: \begingroup \urlstyle{rm}\Url}\fi

\bibitem[Facchinei and Pang(2003)]{facchinei2003}
Francisco Facchinei and Jong-Shi Pang.
\newblock \emph{Finite-Dimensional Variational Inequalities and Complementarity
  Problems}.
\newblock Springer, 2003.
\newblock \doi{10.1007/b97543}.

\bibitem[Korpelevich(1976)]{korpelevich1976}
Galina~M. Korpelevich.
\newblock The extragradient method for finding saddle points and other
  problems.
\newblock \emph{Ekonomika i Matematicheskie Metody}, 12\penalty0 (4):\penalty0
  747--756, 1976.

\bibitem[Rockafellar(1976)]{rockafellar1976}
R.~Tyrrell Rockafellar.
\newblock Monotone operators and the proximal point algorithm.
\newblock \emph{SIAM Journal on Control and Optimization}, 14\penalty0
  (5):\penalty0 877--898, 1976.
\newblock \doi{10.1137/0314056}.

\bibitem[Nemirovski(2004)]{nemirovski2004}
Arkadi Nemirovski.
\newblock Prox-method with rate of convergence {$O(1/t)$} for variational
  inequalities with {Lipschitz} continuous monotone operators and smooth
  convex-concave saddle point problems.
\newblock \emph{SIAM Journal on Optimization}, 15\penalty0 (1):\penalty0
  229--251, 2004.
\newblock \doi{10.1137/S1052623403425629}.

\bibitem[Monteiro and Svaiter(2012)]{monteiro2012}
Renato D.~C. Monteiro and Benar~F. Svaiter.
\newblock Iteration-complexity of a {Newton} proximal extragradient method for
  monotone variational inequalities and inclusion problems.
\newblock \emph{SIAM Journal on Optimization}, 22\penalty0 (3):\penalty0
  914--935, 2012.
\newblock \doi{10.1137/11083085X}.
\newblock URL \url{https://doi.org/10.1137/11083085X}.

\bibitem[Bullins and Lai(2022)]{bullins2020}
Brian Bullins and Kevin~A. Lai.
\newblock Higher-order methods for convex-concave min-max optimization and
  monotone variational inequalities.
\newblock \emph{SIAM Journal on Optimization}, 32\penalty0 (3):\penalty0
  2208--2229, 2022.
\newblock \doi{10.1137/21M1395764}.

\bibitem[Adil et~al.(2022)Adil, Bullins, Jambulapati, and Sachdeva]{adil2022}
Deeksha Adil, Brian Bullins, Arun Jambulapati, and Sushant Sachdeva.
\newblock Optimal methods for higher-order smooth monotone variational
  inequalities.
\newblock arXiv:2205.06167, 2022.
\newblock URL \url{https://arxiv.org/abs/2205.06167}.

\bibitem[Lin and Jordan(2025)]{lin2024}
Tianyi Lin and Michael~I. Jordan.
\newblock {Perseus}: A simple and optimal high-order method for variational
  inequalities.
\newblock \emph{Mathematical Programming}, 209:\penalty0 609--650, 2025.
\newblock \doi{10.1007/s10107-024-02075-2}.

\bibitem[Ostroukhov et~al.(2020)Ostroukhov, Kamalov, Dvurechensky, and
  Gasnikov]{ostroukhov2020}
Petr Ostroukhov, Rinat Kamalov, Pavel Dvurechensky, and Alexander Gasnikov.
\newblock Tensor methods for strongly convex strongly concave saddle point
  problems and strongly monotone variational inequalities.
\newblock arXiv:2012.15595v3, 2020.
\newblock URL \url{https://arxiv.org/abs/2012.15595}.

\bibitem[Chen et~al.(2026{\natexlab{a}})Chen, Zhang, Wang, Liu, Chen, and
  Zhang]{chen2026halpern}
Lesi Chen, Xinliang Zhang, Hengyu Wang, Chengchang Liu, Yongchao Chen, and
  Jingzhao Zhang.
\newblock Halpern iteration accelerates solving high-order monotone variational
  inequalities.
\newblock arXiv:2608.08463v2, 2026{\natexlab{a}}.
\newblock URL \url{https://arxiv.org/abs/2608.08463}.

\bibitem[Zhang et~al.(2026{\natexlab{a}})Zhang, Wu, Zhang, Fang, Li, and
  Lin]{zhang2026}
Haihan Zhang, Wendao Wu, Chenheng Zhang, Cong Fang, Haoxuan Li, and Zhouchen
  Lin.
\newblock Matching higher-order oracle complexity for smooth monotone
  variational inequalities.
\newblock arXiv:2609.13462v2, 2026{\natexlab{a}}.
\newblock URL \url{https://arxiv.org/abs/2609.13462}.

\bibitem[Zhang et~al.(2026{\natexlab{b}})Zhang, Chen, Pan, Liu, Yang, and
  Zhang]{zhangchen2026}
Xinliang Zhang, Lesi Chen, Linxuan Pan, Chengchang Liu, Junchi Yang, and
  Jingzhao Zhang.
\newblock Optimal high-order methods for solving monotone variational
  inequalities.
\newblock arXiv:2609.23557v1, 2026{\natexlab{b}}.
\newblock URL \url{https://arxiv.org/abs/2609.23557}.

\bibitem[Jiang and Yoon(2026)]{jiangyoon2026}
Ruichen Jiang and TaeHo Yoon.
\newblock Anchored extra-proximal methods: Optimal higher-order methods for
  monotone inclusion problems.
\newblock arXiv:2609.30212v2, 2026.
\newblock URL \url{https://arxiv.org/abs/2609.30212}.

\bibitem[Nesterov(2007)]{nesterov2007}
Yurii Nesterov.
\newblock Dual extrapolation and its applications to solving variational
  inequalities and related problems.
\newblock \emph{Mathematical Programming}, 109\penalty0 (2--3):\penalty0
  319--344, 2007.
\newblock \doi{10.1007/s10107-006-0034-z}.

\bibitem[Diakonikolas(2020)]{diakonikolas2020}
Jelena Diakonikolas.
\newblock Halpern iteration for near-optimal and parameter-free monotone
  inclusion and strong solutions to variational inequalities.
\newblock In \emph{Proceedings of the Thirty-Third Conference on Learning
  Theory}, volume 125 of \emph{Proceedings of Machine Learning Research}, pages
  1428--1451, 2020.
\newblock URL \url{https://proceedings.mlr.press/v125/diakonikolas20a.html}.

\bibitem[Cai et~al.(2022)Cai, Oikonomou, and Zheng]{cai2022}
Yang Cai, Argyris Oikonomou, and Weiqiang Zheng.
\newblock Finite-time last-iterate convergence for learning in multi-player
  games.
\newblock In \emph{Advances in Neural Information Processing Systems},
  volume~35, 2022.
\newblock URL
  \url{https://proceedings.neurips.cc/paper_files/paper/2022/hash/db2d2001f63e83214b08948b459f69f0-Abstract-Conference.html}.

\bibitem[Chen et~al.(2025)Chen, Liu, Luo, and Zhang]{chen2025}
Lesi Chen, Chengchang Liu, Luo Luo, and Jingzhao Zhang.
\newblock Solving convex-concave problems with {$\mathcal{O}(\epsilon^{-4/7})$}
  second-order oracle complexity.
\newblock In \emph{Proceedings of the Thirty Eighth Conference on Learning
  Theory}, volume 291 of \emph{Proceedings of Machine Learning Research}, pages
  952--982, 2025.
\newblock URL \url{https://proceedings.mlr.press/v291/chen25a.html}.

\bibitem[Chen et~al.(2026{\natexlab{b}})Chen, Zhang, Liu, Li, Luo, and
  Zhang]{chen2026saddle}
Lesi Chen, Xinliang Zhang, Chengchang Liu, Junru Li, Luo Luo, and Jingzhao
  Zhang.
\newblock Solving convex-concave problems with $\widetilde{\mathcal
  o}(\epsilon^{-4/(3p+1)})$ $p$th-order oracle complexity.
\newblock arXiv:2604.19462v1, 2026{\natexlab{b}}.
\newblock URL \url{https://arxiv.org/abs/2604.19462v1}.

\bibitem[Jang and Ryu(2026)]{jang2026}
Uijeong Jang and Ernest~K. Ryu.
\newblock Higher-order derivatives do not accelerate the computation of fixed
  points.
\newblock arXiv:2607.05947v1, 2026.
\newblock URL \url{https://arxiv.org/abs/2607.05947v1}.

\end{thebibliography}

\endgroup
\end{document}
